\problemname{Planet X}
The year is 2109 and a group of scientists has just discovered ``Planet X'',
a previously unknown planet here in our own solar system,
beyond Pluto's orbit. The research group immediately sends out
a probe to take measurements, and shortly after they receive measurement data.

The scientists are specifically interested in what the surface of Planet X looks like.
We represent the surface as an $N \times M$ grid, where each cell
has a height between $0$ and $9$.

A measuring instrument on the probe has managed to measure the specific height
of some, but not all, cells. From the chemical composition of the surface, we know that it is not
very steep on the planet: the height between two
adjacent cells (cells that share an edge) can never differ
by more than one.

Now the scientists need your help to extract as much information
from this data as possible. More precisely, they want you, given the heights
of some of the cells, to find the heights of all other cells that can be uniquely determined.

\section*{Input}
The first line contains two integers $1 \le N,M \le 10$,
the number of rows and the number of columns of the grid respectively.
Then follow $N$ lines with $M$ characters each.
The $j$-th character on line $i$ is a \texttt{.} if
no value for this cell exists, and is a digit between
$0$ and $9$ corresponding to the height of the cell otherwise.

At least one cell contains a digit.

\section*{Output}
The program should print $N$ lines with $M$ characters each:
the grid as it looks after
correct heights have been filled in for all cells where the height can be determined.

\section*{Scoring}
Your solution will be tested on a set of test groups.
To earn points for a group, you must pass all test cases in that group.

\noindent
\begin{tabular}{| l | l | l |}
  \hline
  \textbf{Group} & \textbf{Points} & \textbf{Constraints} \\ \hline
  $1$    & $20$        &  $N = 1$\\ \hline
  $2$    & $20$        &  $1 \leq M \cdot N \leq 8$\\ \hline
  $3$    & $60$        &  No additional constraints. \\ \hline
\end{tabular}
